\documentclass[11pt]{article}

\usepackage{ulem}
%
\usepackage{amsmath}
%
\usepackage{amssymb}
%
\usepackage{enumerate}
%
\usepackage[width=7in,height=10in,top=0.5in,centering]{geometry}
%
\usepackage{graphicx}
%
\parindent0pt
\fboxsep=0.25cm
%
\newcommand{\Heading}[2]{\fbox{{\Large\textbf{#1}}}\hfill\fbox{\textbf{#2}}\par\vspace{0.1in}}
%
\newcommand{\Name}[1]{\textbf{Names}:\quad#1\hrulefill}
%
\newcommand{\DueDate}[1]{\textbf{Due Date:}\quad#1\par\medskip}
%
\DeclareMathOperator{\arcsec}{arcsec}
\DeclareMathOperator{\arccot}{arccot}
\DeclareMathOperator{\arccsc}{arccsc}
%
\newcounter{example}[section]
\newcounter{exercise}[section]
\newcommand{\important}[2]{\begin{center}\fbox{\begin{minipage}{.97\textwidth}\noindent\textbf{#1}\par#2\end{minipage}}\end{center}}
\newcommand{\defn}[1]{\begin{center}\fbox{\begin{minipage}{.97\textwidth}\reversemarginpar\marginnote{\textbf{Definition}}#1\end{minipage}}\end{center}}
\newcommand{\example}[0]{\vspace{0.1in}{\textcolor{blue}{\textbf{Example \stepcounter{example}\theexample.}}}\hspace{0.1in}}
\newcommand{\exercise}[0]{\vspace{0.1in}{\textcolor{blue}{\textbf{Exercise \stepcounter{exercise}\Alph{exercise}.}}}\hspace{0.1in}}

%%%%%%%%%%%% BEGIN DOCUMENT %%%%%%%%%%%%%%%
%%%%%%%%%%%% BEGIN DOCUMENT %%%%%%%%%%%%%%%
%%%%%%%%%%%% BEGIN DOCUMENT %%%%%%%%%%%%%%%
%%%%%%%%%%%% BEGIN DOCUMENT %%%%%%%%%%%%%%%
%%%%%%%%%%%% BEGIN DOCUMENT %%%%%%%%%%%%%%%

\begin{document}

{\Large\bf MATH 1190/1210 Exhaustive Proctoring Instructions}

\begin{itemize}
\item 6:15 - 6:30 PM: Pick up tests.
\begin{itemize}
    \item Pick up the test from the coordinator. 
    \item You are highly recommended to bring extra writing utensils in case students need them.
    %\item Bring tape and one or more signs to guide students get their IDs checked as they enter the test room.
\end{itemize}
\item 6:30 - 6:40 PM: Arrive at the test room and set up the podium.
\begin{itemize}
    \item Project timer for the length of the test to the screen.
    \item Project general test instructions. (Or write on the board)
    \item Test the microphone (if there is one).
    %\item If there are already students present, check their IDs and mark their names off.
\end{itemize}

\item 6:40 - 6:55 PM: ID Checking (This part might take longer, which is fine. Don't have to start exactly at 7pm as long as students get their full test length. You can also choose to check IDs during or after the test.)
    \begin{itemize}
    \item Have a proctor stand around the main entrance of the room to check IDs. 
    \item If a student didn't bring an ID, get their name and check them again the photo roster on Canvas (can be done quietly during the test as long as you know who it is)
    \item The other proctor should stay with the tests and guide the students toward the ID checking station. Also use this time to answer any individual student questions, if any.
    %\item Ask students to turn phones off \& put away all personal items except writing utensils \& erasers
    %\item No calculators, notes, dictionaries, or other aids are allowed
    %\item Students may not use their own extra paper for scrap work; no loose papers should be on the ground around them
    %\item Students' bags/purses/etc should be closed so their contents are not visible
    %\item Walk around the room to make sure students have followed these instructions, one row at a time
    
    \end{itemize}
    
\item 6:55 - 7:00 PM: Prepare Test Room, Instructions and Test distribution
    \begin{itemize}
    \item Help students find seats by having them scoot to the inner-most seats in their rows and sit one seat apart from each other.
    \item Reserve left-handed seats (typically the left-most seat on a row, from the perspective of someone sitting on the row) for left-handed students, when possible.
    \item If a student has not cleared their space properly, ask them to clear their space as instructed and wait by their row until they have done so.
    \item Repeat until the test room is prepared.
    \item Before distributing tests, tell the students they may NOT open the test until you let them know they may begin.
    \item Ask students to fill out their information on the front of their test and to read the instructions as you pass out the tests - we want to make sure nobody has an advantage by starting their test earlier than others.
    %\item Distribute the tests to the students by handing $n$ tests to the person at the end of the row - where $n$ is the number of students seated in that row - and have them pass the tests down the row.
    \item As you distribute the tests, remind them repeatedly not to open the tests until you tell them to do so.
    \item Once you have distributed the tests to the students, {\bf read the instructions to them out loud}.
    %\item In particular, remind them they should not put DNE when they can either put $=\infty$ or $=-\infty$ and that they should put their answers in the associated answer boxes, when available.
    %\item Bring to their attentions that the learning goal of each question is at the top of the page; encourage them to ask themselves to what degree they think their answers demonstrate proficiency with that learning goal to guide how much work they should show
    %\item Tell them they should show their student IDs to submit their tests; they can use a driver's license or passport, if needed.
    %\item Have them check their tests to make sure there are 5 problems on 5 pages, encourage them to complete they problems they feel best about first
    %\item Remind them they should show their IDs to submit their tests again.
    \item Ask if there's any questions about the instructions. If not, start the test.
    \end{itemize}
    
\item 7:00 PM - end time: Taking the test.
    \begin{itemize}
    \item Start the timer and announce to students that they may begin.
    \item Give time cues as you see fit. Make sure to give a 5-minute warning. %Gently let them know when they have 15 minutes and 5 minutes remaining.
    \item When giving the 5-minute warning, remind them they should show their student IDs to submit their tests.
    \item Rove the room; watch for devices hidden on laps and notes hidden in sleeves on caps or water bottles; listen for weird sounds that resemble communication; watch eyes, particularly if a student is trying to cover them up.
    \item You may ask suspicious students to move seats.
    %\item You may take pictures or videos of suspicious students with your phone if the offense can be easily recorded.
    \item {\bf The only questions we can answer are ones that clarify what the question is asking}, in an English language sense. So we cannot clarify the meaning of any mathematical notations or definitions.
    \item If a student asks about something about the course content, we generally cannot answer that. I usually say, "unfortunately I can't answer that, use your best judgement".
    \item Our responses to their questions should {\bf never} suggest what they should do to solve a problem; some have been trained in grade school that they can get hints if they just keep asking questions, but hints usually defeat the purposes of the test problems.
    \item If you suspect there's any typos on the test, please do not confirm or deny the typo to any student. Message the proctors on Discord to confirm the typo before making any announcement.
    \end{itemize}
    
\item Submission
    \begin{itemize}
    \item Once time is up, ask students to submit their tests. If a student continues to write, gently remind them time is up and they need to stop writing. If they continue to write, ask them to submit the test first. %(It helps to designate one proctor on this task, since ideally most of the students would be lining up at the podium waiting to submit their test.)
    %\item Students should submit their tests with their IDs (unless you know their names).
    %\item Mark their names off on the shared attendance sheet.
    \item \textbf{Sort the tests in the order of their sections.}
    \item Bring tests back to the department.
    \end{itemize}
    
\item Scanning and Storing
    \begin{itemize}
    \item Scan the tests either on the night of or during the day after.
    \item Upload scans to Gradescope to make sure there's no issues.
    \item Once scanning is complete and uploaded, bring the physical tests back to Josh for record keeping.
\end{itemize}
    
\end{itemize}

\end{document}